% ===================================================================
%  BAGAN No. 7 — Desa pada musim dingin. Rumah tani pada musim semi.
%
%  Traduction, faite en 2026, en indonesien d'aujourd'hui, sur l'ido
%  de texto/io. Regles de la colonne : voir 10-bagan-01.tex. Deux
%  scenes, deux numerotations : les cles portent c1 et c2. Le bloc
%  t07-c1-07-1 tient DEUX alineas sous une seule cle — ordo.py le
%  marque « ¶2 ».
%
%  LE BLOC t07-c2-12-1 PASSE ICI SANS EFFORT, ET C'EST LA MEILLEURE
%  PREUVE DE CE QU'IL COUTAIT AILLEURS. Ce bloc a fait trebucher le
%  coreen, le tamoul et l'ourdou l'un apres l'autre : l'ido y pose
%  « fromaji (64), qui gutifas de la tabulo pozita sur la kuvego
%  (63) », une relative APRES son antecedent, et les trois langues a
%  tete finale devaient soit couper la phrase, soit trouver une
%  relative postposee. L'indonesien a « yang », qui suit son
%  antecedent exactement comme le « qui » de l'ido, et la phrase se
%  traduit mot pour mot. Trois colonnes ont paye ce bloc ; la
%  quatrieme ne le voit meme pas.
%
%  LE MARRON CHAUD (42) OUVRE UN TROISIEME CAS, ET IL FALLAIT UNE
%  CINQUIEME COLONNE POUR LE TROUVER. Jusqu'ici le releve n'avait que
%  deux reponses : le coreen avait le MOT et la SCENE (밤, les
%  chataignes grillees de rue) ; le telougou, le marathe, l'egyptien
%  et le tamoul n'avaient ni l'un ni l'autre et devaient periphraser.
%  L'ourdou avait les deux aussi. L'indonesien a le mot et pas la
%  scene : berangan est le nom malais du chataignier — Castanopsis
%  pousse a Sumatra —, mais personne ne vend de marrons chauds dans
%  une rue de Java, ou il fait trente degres. Le releve a donc trois
%  cas, et non deux : le mot sans la chose, la chose sans le mot, et
%  le mot sans la SCENE.
%
%  ET LA NEIGE EST TOUT ENTIERE UN VOCABULAIRE D'ECOLE, COMME LES
%  SAISONS DU TABLEAU 3. L'Indonesie n'a de neige que sur les cimes
%  de Papouasie. Mais la langue ne l'emprunte pas : salju vient de
%  l'arabe thalj et il est ancien ; le reste se FABRIQUE sur des
%  racines malaises — tempat luncur la glissoire (33), kereta luncur
%  la luge (32), sepatu luncur le patin (26), bola salju la boule de
%  neige (38), patung salju le bonhomme (39). Cinq objets qu'aucun
%  lecteur n'a touches, cinq mots transparents batis sur « luncur »,
%  glisser. Une langue qui n'a pas la chose ne se tait pas
%  forcement : elle peut composer.
%
%  Seul « es » (25), la glace, est neerlandais — de ijs. Il designe
%  ce qu'on trouve dans un verre, et c'est le seul objet de la serie
%  hivernale qu'un Indonesien manie tous les jours.
% ===================================================================

%%K t07-noto-1 noto
\VUnotes{143.2mm}{%
(*) Untuk perincian tentang makanan, lihat bagan 6 dan 13.%
}
%%K t07-apar-1 apar
\VUsaut{4.0mm}
\VUcentre{13.2pt}{90}{SERI KEDUA}
\VUsaut{3.0mm}
\VUfilet{16mm}
\VUsaut{5.6mm}
\VUtitre{15.0pt}{60}{BAGAN No. 7}
\VUsaut{3.0mm}

%%K t07-tit-1 sub
\VUcentre{12.0pt}{0}{\textit{Adegan Pertama.}}
\VUsaut{3.0mm}

%%K t07-tit-2 sub
\VUpk{10.2pt}{40}{Desa pada musim dingin.}
\VUsaut{4.0mm}

%%K t07-c1-01-1 p
1. --- Pada liburan tahun baru saya menemani ayah saya ke Novurbo,
sebuah desa kecil tempat ia harus mengurus beberapa urusan dagang.

%%K t07-c1-02-1 p
2. --- Kami naik \VUgras{kereta pos} \textsuperscript{(11)} yang berjalan
antara kota dan kampung itu ; tetapi karena udara sangat dingin,
perjalanan dengan kendaraan yang kurang nyaman itu tidak begitu
menyenangkan.

%%K t07-c1-03-1 p
3. --- Karena itu kami merasa sungguh senang ketika kami melihat
\VUgras{kusir} menghentikan keretanya di alun-alun, di depan
\VUgras{penginapan} \textsuperscript{(2)} \textit{Beruang Putih}.
\VUgras{Pemilik penginapan} \textsuperscript{(4)} sedang menunggu kami
di ambang pintu rumahnya sambil dengan tenang mengisap
\VUgras{pipanya.}

%%K t07-c1-03-2 p
Ia mengajak kami masuk, dan saya tidak menunggu diminta dua kali untuk
duduk di depan api ranting yang berderak di tungku. Ketika itu saya
mengejek habis-habisan \VUgras{angin utara} yang tajam, yang
mendecitkan \VUgras{papan nama} \textsuperscript{(3)} tua itu dan
membawa pergi dalam pusaran hitam \VUgras{asap}
\textsuperscript{(1)} yang keluar dari cerobong.

%%K t07-c1-04-1 p
4. --- Waktu makan siang telah tiba. Tanpa menunggu lebih lama, kami
makan dengan lahap santapan yang sederhana tetapi lezat yang
dihidangkan kepada kami (*).

%%K t07-tit-3 sub
\VUpk{10.2pt}{40}{Beberapa kunjungan.}
\VUsaut{3.0mm}

%%K t07-c1-05-1 p
5. --- Sesudah makan siang saya menemani ayah saya, yang menjadi agen
asuransi, mengunjungi para pelanggannya. Mula-mula kami mengunjungi
\VUgras{pandai besi} \textsuperscript{(5)}, yang tinggal dekat
penginapan. Ia sedang sibuk memasang ladam pada seekor kuda. Saya
mengagumi kekuatan pembantunya, yang memegang teguh \VUgras{kuku kuda}
itu di atas celemek kulitnya, sementara pandai besi memasang
\VUgras{ladam} \textsuperscript{(6)} dengan pukulan palu yang kuat.

%%K t07-c1-06-1 p
6. --- Kemudian kami pergi ke \VUgras{tukang sepatu}
\textsuperscript{(7)} desa itu, si tua Iakobus. Meskipun papan namanya
sebuah \VUgras{sepatu bot} yang sangat indah, ia puas dengan memasang
sol baru pada sepasang sepatu tua yang berlubang.

%%K t07-c1-06-2 p
Orang tua itu berkata kepada kami sambil tertawa : « Di kota saya
adalah « pembuat sepatu » dan saya tidak punya pelanggan. Di sini
saya « tukang perbaiki sepatu » dan pekerjaan banyak sekali. »

%%K t07-c1-07-1 p
7. --- Ia bercerita kepada kami tentang tetangganya, \VUgras{tukang
cukur} \textsuperscript{(8)}, yang pelanggannya tidak puas. \VUgras{Pisau
cukurnya} selalu sumbing dan \VUgras{guntingnya} tidak pernah memotong
dengan baik.

Tetapi karena ia satu-satunya tukang cukur di desa itu, tentu saja
orang terpaksa bercukur dan berpangkas rambut kepadanya. Lagi pula ia
lelaki yang kikir dan tidak ramah.

%%K t07-c1-08-1 p
8. --- Untunglah \VUgras{tetangga} yang lain tidak seperti dia.
Misalnya \VUgras{tukang kereta} \textsuperscript{(9)} adalah lelaki yang
baik hati, rajin dan suka menolong. Menyenangkan sekali memperbaiki
kereta kepadanya. Pekerjaannya selalu tuntas.

%%K t07-c1-09-1 p
9. --- Si tua Iakobus juga tidak mengeluh tentang \VUgras{tukang
pelana} \textsuperscript{(10)}, yang kadang-kadang ia bantu menjahit
\VUgras{tali kekang} bila pekerjaan mendesak.

%%K t07-c1-09-2 p
Ayah saya menanyakan keluarganya : « Semua sehat. Cucu perempuan saya
ada di \VUgras{sekolah} \textsuperscript{(12)}. \VUgras{Guru
perempuannya} \textsuperscript{(13)} sangat senang kepadanya ; ia
\VUgras{murid} \textsuperscript{(14)} yang baik. Ia tidak seperti anak
lelaki saya yang nakal ; ia hanya memikirkan bermain. \VUgras{Guru}
\textsuperscript{(15)} tidak berhasil mengajarkan apa pun kepadanya. »

%%K t07-tit-4 sub
\VUsaut{3.0mm}
\VUpk{10.2pt}{40}{Obrolan desa.}
\VUsaut{3.0mm}

%%K t07-c1-10-1 p
10. --- Orang tua itu kemudian menceritakan kepada kami berita-berita
desa. Sebuah sekolah baru akan dibangun, sebab yang berada di
\VUgras{balai desa} sudah menjadi terlalu kecil. Itu mudah, sebab
cukup membeli sebagian kebun \VUgras{tukang giling.} Hal itu akan
sangat menguntungkan lelaki yang malang itu. Karena udara dingin,
\VUgras{batu-batu giling} pada \VUgras{penggilingan}
\textsuperscript{(16)} tidak dapat lagi menggiling gandum, sebab
\VUgras{rodanya} \textsuperscript{(17)} dihentikan beku oleh \VUgras{es}
\textsuperscript{(25)} \VUgras{sungai kecil} \textsuperscript{(23)}.

%%K t07-c1-11-1 p
11. --- « Tentang bangunan, sudahkah Anda melihat \VUgras{gereja}
\textsuperscript{(18)} kami yang baru? Ia sangat menarik di bawah
selimut saljunya. \VUgras{Burung-burung gagak} \textsuperscript{(19)},
yang tidak lagi mengenali menara loncengnya yang lama, hinggap dengan
sedih di pohon besar \VUgras{kuburan} \textsuperscript{(20)}. »

%%K t07-c1-12-1 p
12. --- Pikiran tentang kuburan itu mengingatkan tukang sepatu pada
orang-orang yang dikenal ayah saya dan yang meninggal tahun lalu. «
Berapa banyak \VUgras{makam} \textsuperscript{(21)}, tuan yang baik,
bertambah pada yang lain, di bawah \VUgras{pohon sipres}
\textsuperscript{(22)}! Maut telah menebas notaris tua kami. Seluruh
penduduk desa menghadiri \VUgras{pemakamannya.} »

%%K t07-c1-13-1 p
13. --- « Itulah penggantinya, yang sedang berseluncur di sungai
kecil. Ia baru saja datang untuk memperbaiki tali \VUgras{sepatu
luncurnya} \textsuperscript{(26)} yang putus. »

%%K t07-c1-14-1 p
14. --- « Itu juga, di tepi sungai kecil, \VUgras{penjaga desa}
\textsuperscript{(27)} yang baru. Ia bekas polisi tentara dan menakuti
anak-anak nakal desa. Mereka langsung lari begitu melihat
\VUgras{pembalut betisnya} \textsuperscript{(29)} dan \VUgras{topi
dinasnya} \textsuperscript{(28)}. »

%%K t07-c1-15-1 p
15. --- « Ha! itu si tua Iosef yang menghela \VUgras{gerobak kecil
berisi ikatan kayu} \textsuperscript{(30)} untuk tukang roti. ---
Benar, kata ayah saya ; saya mengenali pasangan \VUgras{bagalnya}
\textsuperscript{(31)}. Saya memang perlu berbicara dengannya ; saya
akan memanfaatkan kesempatan ini. Sampai jumpa, pak Iakobus. »

%%K t07-c1-16-1 p
16. --- Tepat ketika kami keluar, anak-anak desa meninggalkan sekolah
dan berlari menghambur ke \VUgras{tempat luncur} \textsuperscript{(33)}
dengan risiko \VUgras{jatuh} \textsuperscript{(34)} dan mematahkan
anggota badan. Salah seorang dari mereka mengambil \VUgras{kereta
luncurnya} \textsuperscript{(32)} dari tukang kereta, mendudukkan
seorang kawannya di atasnya, dan berangkat sambil berlari.

%%K t07-c1-17-1 p
17. --- Yang lain menghambur ke \VUgras{onggokan salju}
\textsuperscript{(37)} di dekat \VUgras{jembatan} \textsuperscript{(40)}
dan mulai menyerang \VUgras{patung salju} \textsuperscript{(39)} yang
mereka dirikan kemarin dan yang mereka pakaikan topi, pipa dan tas
sekolah tersandang.

%%K t07-c1-18-1 p
18. --- Mereka tidak begitu peduli pada angin yang membekukan dan
pada dinginnya. Sebagian besar bahkan tidak memakai \VUgras{syal}
\textsuperscript{(35)}, dan untuk menghangatkan diri kembali sesudah
perang \VUgras{bola salju} \textsuperscript{(38)}, cukuplah bagi mereka
segenggam \VUgras{buah berangan} \textsuperscript{(42)} yang sangat
panas, yang dibeli dari \VUgras{penjual perempuan} tua Jakobina, yang
menggigil kedinginan di bawah \VUgras{selendangnya}
\textsuperscript{(43)} yang terlalu tipis.

%%K t07-tit-5 sub
\VUsaut{3.0mm}
\VUcentre{12.0pt}{0}{\textit{Adegan Kedua.}}
\VUsaut{3.0mm}

%%K t07-tit-6 sub
\VUpk{10.2pt}{40}{Rumah tani pada musim semi.}
\VUsaut{4.0mm}

%%K t07-c2-01-1 p
1. --- Meskipun ada segala kesenangan musim dingin, kami menyambut
kembalinya \VUgras{musim semi} dengan kegembiraan yang dalam.

%%K t07-c2-01-2 p
Alangkah riangnya pemandangan halaman rumah tani pada sore hari, di
bulan Mei!

%%K t07-c2-02-1 p
2. --- Di kaki langit \VUgras{matahari} \textsuperscript{(1)} perlahan
terbenam, membanjiri \VUgras{pedesaan} \textsuperscript{(3)} dengan
\VUgras{sinarnya} \textsuperscript{(2)} yang ungu. \VUgras{Pembajak}
\textsuperscript{(6)} yang rajin bergegas membajak \VUgras{ladangnya}
\textsuperscript{(4)}.

%%K t07-c2-02-2 p
Dengan \VUgras{bajaknya} \textsuperscript{(5)}, yang ditarik seekor
\VUgras{kuda} \textsuperscript{(7)} yang kuat, ia menarik \VUgras{alur}
\textsuperscript{(8)}, tempat \VUgras{penabur benih}
\textsuperscript{(9)} akan melemparkan \VUgras{benih} segenggam demi
segenggam, yang akan menghasilkan bulir-bulir keemasan.

%%K t07-c2-03-1 p
3. --- \VUgras{Para pemotong rumput} \textsuperscript{(11)}, dengan
sabit mereka, telah memotong rumput padang ; para penjemur telah
mengeringkan \VUgras{jerami kering} \textsuperscript{(10)} dengan
membalik-baliknya memakai \VUgras{garpu rumput} \textsuperscript{(12)}
dan garu, dan inilah mereka kembali.

%%K t07-c2-03-2 p
Sebagian berjalan di depan gerobak berat yang ditarik sapi ; mereka
memikul perkakas mereka di bahu. Sebagian lain, yang lebih malas,
duduk dengan nyaman di atas jerami kering yang harum.

%%K t07-c2-04-1 p
4. --- Sambil lewat mereka memberi salam ramah kepada \VUgras{tukang
kebun} \textsuperscript{(16)}, yang sibuk di \VUgras{kebun buah}
\textsuperscript{(13)} memangkas \VUgras{pohon-pohon buah} dan
mengokulasi \VUgras{pohon pir} \textsuperscript{(14)}. \VUgras{Pohon
prem} \textsuperscript{(15)} mulai berbunga, dan di \VUgras{kebun
sayur} \textsuperscript{(17)} yang dipupuk dengan baik, sayuran, kubis
dan selada sudah tumbuh bagus.

%%K t07-c2-04-2 p
Di atas tembok kecil, seekor \VUgras{merak} \textsuperscript{(18)} yang
indah membentangkan ekornya, supaya bulunya yang sangat indah
dikagumi.

%%K t07-tit-7 sub
\VUpk{10.2pt}{40}{Halaman rumah tani.}
\VUsaut{3.0mm}

%%K t07-c2-05-1 p
5. --- Pintu \VUgras{kandang kuda} \textsuperscript{(19)} terbuka, dan
orang melihat palungan dan \VUgras{alas jerami} \textsuperscript{(23)}
\VUgras{kuda betina} \textsuperscript{(21)} yang baru saja dilepas
\VUgras{tukang kuda} \textsuperscript{(20)}. \VUgras{Anak kuda}
\textsuperscript{(22)}, yang tampak begitu lucu dengan kakinya yang
panjang, mengikuti induknya sambil melonjak-lonjak.

%%K t07-c2-06-1 p
6. --- Di dekat \VUgras{sumur} \textsuperscript{(29)}, \VUgras{sapi-sapi
betina} \textsuperscript{(25)} pergi minum di \VUgras{palungan}
\textsuperscript{(26)} kayu yang menjadi tempat minum mereka.
\VUgras{Anak sapi} \textsuperscript{(24)} memanfaatkan kesempatan itu
untuk menyusu, dan \VUgras{sapi jantan} \textsuperscript{(27)}, yang
mencium bau harum lumbung jerami, melenguh gembira dan mengangkat
kepalanya dengan bangga bersama \VUgras{tanduknya}
\textsuperscript{(28)} yang kuat.

%%K t07-c2-07-1 p
7. --- Semua orang bekerja. \VUgras{Pelayan perempuan}
\textsuperscript{(32)} memegang \VUgras{tali} \textsuperscript{(31)}
sumur dengan tangan. Ia menimba air dengan \VUgras{ember}
\textsuperscript{(30)} untuk memberi minum ternak. Di dekat
\VUgras{lumbung} \textsuperscript{(33)}, seorang lelaki menggaru
jerami, dan di depan pintu \VUgras{rumah} \textsuperscript{(34)},
seorang \VUgras{perempuan pemintal} \textsuperscript{(35)} tua memintal
\VUgras{rami} atau \VUgras{henep} dengan kincir dan \VUgras{jantranya,}
seperti pada zaman dahulu.

%%K t07-tit-8 sub
\VUsaut{3.0mm}
\VUpk{10.2pt}{40}{Petani.}
\VUsaut{3.0mm}

%%K t07-c2-08-1 p
8. --- \VUgras{Petani} \textsuperscript{(36)} memimpin semua pekerjaan
itu dan memberi petunjuk kepada para pekerjanya. Ia menyuruh
menyimpan di bawah atap \VUgras{garu} \textsuperscript{(40)} dan
\VUgras{cangkul} \textsuperscript{(39)} yang tertinggal di dekat
\VUgras{kandang} \textsuperscript{(38)} \VUgras{anjing}
\textsuperscript{(37)}.

%%K t07-c2-09-1 p
9. --- Teriakannya menakuti seekor \VUgras{merpati}
\textsuperscript{(41)}, yang terbang dengan sayap terkembang penuh ke
dalam \VUgras{rumah merpati} \textsuperscript{(42)}, dan
\VUgras{ayam-ayam betina} \textsuperscript{(43)} yang bergegas kembali
ke \VUgras{kandang ayam} \textsuperscript{(44)}.

%%K t07-c2-10-1 p
10. --- Di depan \VUgras{kandang domba} \textsuperscript{(48)}, inilah
\VUgras{gembala} \textsuperscript{(45)} tua yang membawa pulang
\VUgras{kawanannya} \textsuperscript{(49)}. Sambil memegang
\VUgras{tongkatnya} \textsuperscript{(46)}, yang tidak pernah lepas
darinya, seperti juga \VUgras{blusnya} \textsuperscript{(47)} yang
telah pudar warnanya, ia menggiring ke kandang \VUgras{domba-domba
jantan} \textsuperscript{(50)}, \VUgras{domba-domba betina}
\textsuperscript{(53)}, \VUgras{anak-anak domba}
\textsuperscript{(54)}, \VUgras{kambing-kambing betina}
\textsuperscript{(51)} dan \VUgras{anak-anak kambing}
\textsuperscript{(52)}. \VUgras{Anjingnya} \textsuperscript{(55)} yang
setia, Argus, datang paling akhir dan mendorong yang tertinggal dengan
gonggongannya.

%%K t07-c2-11-1 p
11. --- Segala suara dan gerak itu tidak mengganggu ketenangan
\VUgras{sapi perah} \textsuperscript{(56)} yang baik, yang
membunyikan \VUgras{lonceng kecilnya} \textsuperscript{(57)} sementara
ia diperah di depan pintu \VUgras{kandang} \textsuperscript{(58)}.

%%K t07-c2-12-1 p
12. Ia tampak mengerti bahwa ia harus memberikan susunya, supaya
\VUgras{istri petani} \textsuperscript{(60)} dapat membuat
\VUgras{mentega} di dalam \VUgras{pengaduk mentega}
\textsuperscript{(61)} atau \VUgras{keju} \textsuperscript{(64)} yang
sangat baik, yang menetes dari papan yang diletakkan di atas
\VUgras{bak besar} \textsuperscript{(63)}.

%%K t07-c2-13-1 p
13. --- Seluruh \VUgras{tempat susu} \textsuperscript{(59)} dijaga
sangat bersih oleh \VUgras{buruh tani} \textsuperscript{(65)}, yang baru
saja mencuci \VUgras{periuk-periuk susu} \textsuperscript{(62)} di
dalam ember besar yang dibawanya di tangan.

%%K t07-tit-9 sub
\VUsaut{3.0mm}
\VUpk{10.2pt}{40}{Kandang unggas.}
\VUsaut{3.0mm}

%%K t07-c2-14-1 p
14. --- Dengan terompah kayunya yang tebal, ia berjalan agak berat,
sehingga membuat lari \VUgras{kalkun jantan} \textsuperscript{(68)},
\VUgras{itik-itik betina} \textsuperscript{(66)}, \VUgras{itik-itik
jantan} \textsuperscript{(67)} dan \VUgras{angsa-angsa}
\textsuperscript{(69)}, yang membuka lebar \VUgras{paruh mereka}
\textsuperscript{(70)} dan berteriak gelisah.

%%K t07-c2-14-2 p
\VUgras{Ayam jantan} \textsuperscript{(71)}, yang lebih suka berkelahi,
menegakkan \VUgras{jenggernya} \textsuperscript{(72)} yang merah,
mengibaskan bulu \VUgras{ekornya} \textsuperscript{(73)}, dan
memperdengarkan « kukuruyuk » yang nyaring.

%%K t07-c2-15-1 p
15. --- \VUgras{Induk ayam} \textsuperscript{(74)} mengais di atas
\VUgras{pupuk kandang} \textsuperscript{(81)} bersama \VUgras{anak-anak
ayamnya} \textsuperscript{(75)}. \VUgras{Anak babi}
\textsuperscript{(78)} lari kepada induknya, \VUgras{induk babi}
\textsuperscript{(77)} yang besar, yang kita lihat di dekat kandang
babi.

%%K t07-c2-16-1 p
16. --- Di dalam kotak-kotaknya, \VUgras{kelinci-kelinci}
\textsuperscript{(79)} makan dengan lahap \VUgras{rumput}
\textsuperscript{(80)} yang lembut, yang dibawakan anak perempuan
petani itu. Ioanina sangat menyayangi kelincinya, dan setiap hari,
sesudah mengambil \VUgras{telur-telur} \textsuperscript{(76)} segar di
kandang ayam, ia datang menengok kawan-kawan kecilnya.

\VUsaut{6.0mm}
\VUfilet{22mm}
